\section{Polynomial cells and the increment theorem}\label{ap:setup}

Fix $k\ge3$ and put $s=k-2$. We seek density increments on regions
that retain every polynomial constraint introduced by an earlier
increment. The regions will be triangular polynomial cells. Their
integer coordinates are determined successively, and their residual
widths are indexed by weight. The central estimate will bound the
relative width loss at each weight using only higher-weight widths.

Every degree and descent depth below is fixed in terms of $s$;
dimensions may grow. A degree bound may be replaced by any fixed
larger one. We work on rectangular integer boxes, with uniform
probability averages and max norms, allowing dimension factors in
norm comparisons. A \emph{slice} is an integer-affinely reindexed
residue subbox. Different coordinate strides are permitted. Its
logarithmic cost includes their logarithms and the logarithms of
inverse relative side lengths. Sufficient lower bounds on the
original sides absorb all endpoint rounding errors. A function is
\emph{progression-free} when its support contains no nonconstant
integer-vector $k$-term progression.

\subsection{Precision conventions}

Write $\mathcal P(P)=(2+P)^C$ and
$\mathcal B(P)=\exp((2+P)^C)$, with $C$ depending only on the fixed
degrees and allowed to increase. A quantity has polynomial
\emph{logarithmic} cost when its size and, when needed, inverse
size are at most $\mathcal B(P)$. Dimensions, numbers of structural
variables, and polynomial degrees are not measured logarithmically.
Only a bounded number of successive polynomial changes of budget
is implicit; a polynomial number of multiplicative factors is
permitted because their logarithmic costs add. The height of rational
data bounds both numerators and positive denominators; a denominator
bound alone leaves numerators unrestricted.
Multiplying a density threshold by a fixed positive constant changes
its logarithmic budget by only a fixed additive amount. We absorb
such changes in the displayed polynomial bounds, including the fresh
rank bound $d_0$, whenever a discounted threshold is used.

\subsection{Cells and certificates}

Choose a fixed $D\ge s$, whose existence is verified by the weight
argument in Section~\ref{ap:extraction}. Give the spatial variable
$u$ weight one and integer blocks $b_h\in\Z^{d_h}$ weight $h$,
for $1\le h\le D$. The weighted degree of a monomial is the sum
of its variable weights, counted with multiplicity. A triangular cell
is specified by polynomials
\[
 C_h(u,b_{<h})\in\R^{d_h},\qquad \deg_{\rm wt}C_h\le h,
\]
bounded centers $l_h$, and widths $0<w_h\le1/32$. There is no
bound on the coefficients of $C_h$. Integer translations keep the
centers bounded without changing this convention. Zero-dimensional
blocks are allowed. Put
\[
                         Q_h=\log(2/w_h).
\]
Thus $Q_h$ measures the logarithmic precision of the residual at weight~$h$.

For a small $\gamma>0$, let $B^-$ and $B^+$ be the indicators that
there is an integer tuple satisfying respectively
\begin{equation}\label{ap:cell-boxes}
 \|b_h-C_h(u,b_{<h})-l_h\|_\infty\le w_h,
 \qquad
 \|b_h-C_h(u,b_{<h})-l_h\|_\infty\le(1+\gamma)w_h
 \quad(1\le h\le D).
\end{equation}
Triangularity and the narrow widths imply uniqueness of the tuple,
also on slightly padded boxes. Boundary conventions may be chosen
consistently and do not require any regularity of the point
distribution. A \emph{density certificate} is
\begin{equation}\label{ap:certificate}
                 \E fB^- > a\E B^+,
 \qquad 0\le f\le1,
\end{equation}
where $f$ is progression-free. Its excess is not assumed to have a
quantitative lower bound.

We may replace the two indicators temporarily by product smooth
cutoffs: the small cutoff is below the strict indicator, and the
large cutoff is one on the enlarged box and supported slightly
outside it. Their plateaus and transitions are proportional to the
width and gap. When translating small boxes, the ratio of enlarged
to strict coverage volumes is $(1+O(\gamma))^{O(\sum_h d_h)}$.
We choose $\gamma$ from dimension forecasts, independently of all
widths, and pay this ratio by a small multiplicative discount of
the target. Thus no hidden boundary regularity assertion is needed.

\subsection{The triangular increment}

The new cell will retain the old determining coordinates and may add
two kinds of integer coordinates. Fresh \emph{absolute slots} come
from the analytic increment. Preparation may also copy a higher-weight
block into a lower-weight block. The following theorem controls both
additions and the loss of residual width.

\begin{theorem}[Triangular increment]\label{ap:increment}
There are fixed $D=D_k$, $\eta=\eta_k>0$, and polynomials with
degrees depending only on $k$ having the following property.
Suppose $p\ge2$, $a\ge e^{-p}$, and a progression-free input on a
positive-dimensional integer box has a certificate
\eqref{ap:certificate}, with width logs $Q_h=\log(2/w_h)$.
If every root side is sufficiently large with logarithmic cutoff
polynomial in $p$, the dimension data, and all $Q_h$, then either
the hypotheses are impossible or a root slice carries a new
certificate at threshold $(1+\eta)a$.

The new cell retains the determining data through exact integer
reparametrizations and adds at most $d_0=(2+p)^{C_k}$ fresh absolute
slots together with downward preparation copies. Its dimensions
$d'_i$ and width logs $Q'_i$ satisfy
\begin{align}
 d'_i-d_i
   &\le d_0+P_i(p,d_{i+1},\ldots,d_D),\label{ap:dimension-rule}\\
 Q'_i-Q_i
   &\le P'_i(p,\text{all dimension forecasts},Q_{i+1},\ldots,Q_D).
   \label{ap:width-rule}
\end{align}
The root slice cost has polynomial logarithmic bound in all the
same data, now including every $Q_i$. Polynomial exponents and the
fixed gain are independent of the number of increments. Gap and
discount choices can be made in advance from the dimension and
cut-count forecasts.
\end{theorem}

The main point is the absence of $Q_i,Q_{<i}$ from the right-hand side
of \eqref{ap:width-rule}. For the interval application, reserve
$T=O_k(p)$ increments, starting with no slots and fixed width data.
The dimension recurrence first bounds the top block, then each lower
block. Proposition~\ref{ap:iteration:dimensions} verifies this
recurrence for the actual preparation copies and bounds all
within-round structural counts. It supplies one polynomial forecast
$d_*(p)$ for the entire run, before any numerical precision is chosen.

With that forecast fixed, choose nondecreasing polynomials $R_i$
dominating the width losses in \eqref{ap:width-rule}. Define
\[
 V_D=Q_D(0)+T R_D(p,d_*),\qquad
 V_i=Q_i(0)+T R_i(p,d_*,V_{i+1},\ldots,V_D)
 \quad(i=D-1,\ldots,1).
\]
Descending induction gives $Q_i(t)\le V_i$ for all $t\le T$.
Since $D$ is fixed, every $V_i$ is polynomial in $p$. The logarithmic
root-side requirements and the sum of all logarithmic slice costs are therefore
polynomial as well. Meanwhile the threshold rises by a fixed factor
at every increment and, after $O_k(p)$ steps, exceeds one. A
certificate at threshold at least one is impossible because
$fB^-\le B^+$ pointwise. Section~\ref{ap:iteration} makes this
deduction precise and proves the required dimension and width
recurrences.

It remains to construct the increment with these particular
dependencies. A bound polynomial in all current width logs would
not suffice: repeated substitution could raise the polynomial degree
at every round.

\subsection{Polynomial patches}\label{ap:analytic-objects}

The analytic increment produces a smooth test rather than a cell
indicator. These tests are polynomial patches; their integer slots
will supply the new absolute coordinates. Give integer slots $b_1,\ldots,b_d$ nondecreasing
weights $1\le\omega_i\le q$, and give each spatial parameter weight
one. A degree-$q$ \emph{patch} is
\begin{equation}\label{ap:patch-definition}
 B(u)=\sum_{b\in\Z^d}\Phi(S_u(b)),\qquad
 S_u(b)_i=b_i-A_i(u,b_1,\ldots,b_{i-1}),
\end{equation}
where $\deg_{\rm wt}A_i\le\omega_i$ and
$\Phi:\R^d\to[0,1]$ is Lipschitz and supported in
$[-1/4,1/4]^d$. Its \emph{rank} is the number $d$ of integer slots.
The number of parameters, $d$, and $\log(2+\Lip\Phi)$ enter its
logarithmic complexity; the coefficients of the $A_i$ do not.
Rank zero includes the constant patch $1$.
Successively determining the integer slots shows that at most one
summand is nonzero, so $0\le B\le1$.
Figure~\ref{fig:patch-coordinates} illustrates the distinction between
integer determining slots and small real residuals.
For a nonnegative test $B$, its \emph{score at target $a$} is
$\E(f-a)B$. Positivity means that the mean of $f$ under the
probability weight $B/\E B$ exceeds $a$; a quantitative lower
bound on the score also controls the available mass.
\begin{figure}[!b]
\centering
\begin{tikzpicture}[x=1.9cm,y=1.9cm,>=Stealth,
  every node/.style={font=\small},line cap=round,line join=round]
  \draw[->] (-.18,0)--(2.25,0) node[right] {$b_1$};
  \draw[->] (0,-.18)--(0,2.25) node[above] {$b_2$};
  \foreach \i in {0,1,2} {\foreach \j in {0,1,2} {
    \fill[black!45] (\i,\j) circle (1.25pt);
  }}
  \draw[blue!65!black,fill=blue!12,thick]
    (.75,.625)--(1.25,.875)--(1.25,1.375)--(.75,1.125)--cycle;
  \fill[black] (1,1) circle (1.8pt);
  \draw[thin] (1.03,1.01)--(1.63,1.30);
  \node[anchor=west] at (1.68,1.30) {$(1,1)$};
  \node[below] at (1,0) {$1$};
  \node[left] at (0,1) {$1$};
  \node[below] at (1,-.40) {Integer slots};

  \draw[->,thick] (2.75,1.1)--(3.65,1.1)
    node[midway,above=4pt] {$S_1$};

  \begin{scope}[shift={(4.65,1)}]
    \draw[->] (-.78,0)--(.83,0) node[right] {$r_1$};
    \draw[->] (0,-.78)--(0,.83) node[above] {$r_2$};
    \draw[blue!65!black,fill=blue!12,thick]
      (-.25,-.25) rectangle (.25,.25);
    \fill[black] (0,0) circle (1.8pt);
    \node[below] at (0,-1.40) {Residual coordinates};
  \end{scope}
\end{tikzpicture}
\caption{A degree-two patch of rank two, with slot weights $1,2$, at $u=1$.
Here $A_1(u)=u$, $A_2(u,b_1)=(u^2+ub_1)/2$, and
$S_1(b_1,b_2)=(b_1-1,b_2-(1+b_1)/2)$.
The shaded parallelogram is the inverse image of
$[-1/4,1/4]^2$, the permitted support of the kernel.
Only the integer pair $(1,1)$ lies in it.
For a general patch, at most one integer tuple satisfies the support
constraints, by successive determination of the slots.}
\label{fig:patch-coordinates}
\end{figure}


\subsection{The absolute starting rule}

Positive progression counting first gives an increment on boxes of a
fixed dimension. Relative lifting transfers that rule to arbitrary
dimension while keeping the number of new slots separate from the
other complexity costs. The following form is used at the terminal
boxes of the cell argument.

\begin{lemma}[Dimension-independent fresh rank]\label{ap:absolute}
There are $\xi>0$ and $C=C_k$ with the following property.
Let $p\ge2$, and let $f$ be a progression-free $[0,1]$-valued
function of mean at least $a\ge e^{-p}$ on a box of positive
dimension $n$. Suppose every side is sufficiently large, with
a sufficient logarithmic cutoff polynomial in $2+p+n$. Then,
unless the hypotheses are impossible, $f$ admits a degree-$s$
polynomial patch on a slice with positive score at target $(1+\xi)a$. The number of slots is at most
$d_0=(2+p)^C$, independent of $n$. The logarithms of inverse score,
Lipschitz norm, slice cost, and sufficient side cutoff are
polynomial in $2+p+n$.
\end{lemma}

\begin{proof}
Fix the base dimension $n_0(s)$ and prime-side comparability
constant in Proposition~\ref{counting:relative-target}.
Proposition~\ref{counting:absolute}, Lemma~\ref{counting:conversion},
and Theorem~\ref{shift:required} give an absolute rule on these
fixed-dimensional boxes, with gain $1+\xi_0$ and rank
$(2+p)^{C_k}$. The side cutoff absorbs all three requirements;
the prime sides can be odd.

Choose a small fixed $c>0$. At any actual mean
$\alpha\ge(1-c)a$, the absolute rule has logarithmic budget
$O(1+p)$ and gives positive score at the universal smaller target
$\Lambda=(1+\xi_0)(1-c)a$. If $\Lambda\ge1$, its nonnegative
test gives an impossibility rule instead. In the relative theorem
use the constant old test $1$, of rank zero, at threshold
$(1-c)a$. Its old score is at least $ca$.
Enlarge the resource parameter to a polynomial in $p+n$ to cover
the dimension, absolute cutoff, and $\log(1/(ca))$.
Theorem~\ref{lifting:relative-theorem} proves
Proposition~\ref{counting:relative-target}, including its separate
impossibility clause. Its rank conclusion is $d+d_0$, so the
enlarged resource parameter does not enlarge the fresh rank $d_0$.
Choose its fixed target discount $\epsilon_0$ and then $c$ so that
$(1-\epsilon_0)(1+\xi_0)(1-c)>1+\xi$ for a fixed $\xi>0$.
This gives the claimed rule and all its quantitative bounds.
\end{proof}

\subsection{Scalar rank and scale conventions}\label{ap:scalar-conventions}

Preparation tests the homogeneous parts of the determining polynomials.
For a box with side lengths $H=(H_1,\ldots,H_n)$, write
$H^\nu=\prod_iH_i^{\nu_i}$. A homogeneous scalar or vector
polynomial is \emph{slow} at budget $p$ if its coefficient of
$u^\nu$ has magnitude at most $\mathcal B(p)H^{-\nu}$.
It is \emph{rational} if its coefficients have a common denominator
at most $\mathcal B(p)$, with no numerator bound. The homogeneous
part at the assigned degree or weight is called its \emph{symbol}.
A polynomial is \emph{fully slow} when the same coefficient bounds
hold at every degree, including the constant coefficient, after the
box has been translated to the origin. Thus real estimates use
normalized variables, while denominator estimates use the raw
integer variables. Group-valued symbols are introduced with the
ascent representation in Section~\ref{ap:ascent}.

\subsubsection{Formal rank}\label{ap:rank-conventions}
Suppose a degree-$j$ polynomial $C_j$ takes values in a rational
space $W_j\subseteq\R^{d_j}$ on a box of sides $H$.
It has rank at least $R$ at this scale if every integer row
$\lambda\in\Z^{d_j}$ with $\|\lambda\|\le R$ and
$\lambda|_{W_j}\ne0$ has the following property: the degree-$j$
homogeneous part of $\lambda C_j$ cannot be written as a polynomial
whose coefficient of $u^\nu$ is bounded by $R/H^\nu$ plus a
rational polynomial with common denominator at most $R$.
For a system, this condition is imposed at every layer.
This is the rank condition of Definition~\ref{sampling:rank};
it is different from the slot count called the rank of a patch.

\subsection{Descending precision budgets}

\begin{definition}[Descending budgets]\label{ap:budgets}
At layer $j$, the \emph{warm data} consist of $p$, forecasts for all
dimensions and structural variable counts, gap parameters depending
only on these data, and the widths and numerical budgets already
chosen at layers strictly above $j$. A warm logarithmic bound is
polynomial in these data and does not use any $Q_i$ with $i\le j$.

The \emph{cold preliminary data} may additionally use the current
width log $Q_j$, including its forecast preparation losses, but not
the eventual rank threshold, long path scale, height of the prepared
value space, or parent-side cutoff at layer $j$. The \emph{late data}
are the perturbation, range, rank, and side choices made after all
preliminary modeling and algebraic requirements at that layer.
Their logarithms are polynomial in the cold preliminary budget.
\end{definition}

The descending choice order is opposite to the order in which paths
are sampled through the layers. Auxiliary models may have cold
complexity. Their eventual return to the original variables must
nevertheless have a warm \emph{relative width loss}; the extraction
argument preserves the defining polynomial maps exactly.
Figure~\ref{ap:budget-figure} displays the permitted dependence.

\begin{figure}[t]
\centering
\begin{tikzpicture}[
 >=Stealth,
 box/.style={draw,align=center,text width=4.8cm,minimum height=1.05cm,
             inner sep=5pt,font=\small},
 every edge/.style={draw,->},node distance=1.05cm and 1.1cm]
\node[box] (warm) {Warm data at $j$\\$p$, dimensions, higher-layer budgets};
\node[box,right=of warm] (loss) {Bound for $Q'_j-Q_j$\\independent of $Q_{\le j}$};
\node[box,below=of warm] (own) {Current width $Q_j$\\and forecast preparation losses};
\node[box,below=of loss] (cold) {Cold and late budget $U_j$\\models, rank, path and side scales};
\node[box,below=of cold] (next) {Warm input at layer $j-1$\\higher data may now include $U_j$};
\path (warm) edge (loss)
      (warm) edge (cold)
      (own) edge (cold)
      (loss) edge (cold)
      (cold) edge (next);
\end{tikzpicture}
\caption{The allowed dependence is triangular. There is no arrow from
the current width or its cold outputs back to its relative width loss.
The full dependency schedule is proved in Section~\ref{ap:iteration}.}
\label{ap:budget-figure}
\end{figure}

\subsection{From a cell to the next certificate}

Preparation removes low-rank relations while preserving the
certificate (Proposition~\ref{ap:major:preparation}). Constrained
paths then carry it to terminal boxes
(Proposition~\ref{ap:scalar:forward}), where
Lemma~\ref{ap:absolute} supplies a quantitative score. The passive
and active return arguments carry that score back through the layers
(Proposition~\ref{ap:passive:replacement} and
Theorem~\ref{ap:active:replacement}). Extraction turns the returned
score into a triangular cell
(Proposition~\ref{ap:macro:extraction}). Finally,
Section~\ref{ap:iteration} verifies the dimension and width bounds
and iterates the increment.

\begin{figure}[htbp]
\centering
\begin{tikzpicture}[
  box/.style={draw,rounded corners,align=center,text width=0.39\linewidth,
              inner sep=6pt,font=\small},
  flow/.style={->,thick},node distance=9mm and 8mm]
\node[box] (prepare) {Prepare the cell\\
  rank cuts preserve the certificate\\
  Proposition~\ref{ap:major:preparation}};
\node[box,right=of prepare] (descend) {Descend along constrained paths\\
  reach terminal boxes\\
  Proposition~\ref{ap:scalar:forward}};
\node[box,below=of descend] (absolute) {Apply the absolute increment\\
  obtain a quantitative score\\
  Lemma~\ref{ap:absolute}};
\node[box,left=of absolute] (return) {Return through the layers\\
  passive and active replacement\\
  Sections~\ref{ap:ascent} and~\ref{ap:active}};
\node[box,below=of return] (extract) {Extract a triangular cell\\
  preserve the required losses\\
  Proposition~\ref{ap:macro:extraction}};
\node[box,right=of extract] (iterate) {Close the budgets and iterate\\
  prove the density bound\\
  Section~\ref{ap:iteration}};
\draw[flow] (prepare) -- (descend);
\draw[flow] (descend) -- (absolute);
\draw[flow] (absolute) -- (return);
\draw[flow] (return) -- (extract);
\draw[flow] (extract) -- (iterate);
\end{tikzpicture}
\caption{The main proof route.  Preparation and extraction preserve the
triangular dependence of the width losses; the absolute increment
controls the number of fresh determining coordinates.}
\label{ap:proof-route}
\end{figure}


During preparation the polynomial $C_j$ takes values exactly in a
rational subspace $W_j\subseteq E_j=\R^{d_j}$, initially the full
space. The rank test defined above and in Definition~\ref{sampling:rank}
requires that no integer
row of norm at most $R_j$, nonzero on $W_j$, has a degree-$j$
symbol that is slow plus rational at that threshold. The test is made
on the current path box at a history that has survived the earlier
tests. The histories arriving at each layer carry the subprobability
law of Section~\ref{ap:path:stopped-tree}; stopping removes mass
without renormalizing the survivors.
Its failure probability may charge $Q_{\ge j}$ but not lower
widths. There is no single root equidistribution test at the finest
precision of all layers.

At a successful layer use a Haar center in
$W_j/(W_j\cap\Z^{d_j})$, conditioned to the feasible range
specified in Section~\ref{ap:path:stopped-tree}. The lift has exact identity
\[
                     C_j-c=y+\beta,
\]
where $y\in\R^{d_j}$ is the small real residual and
$\beta\in\Z^{d_j}$ is the integer lift. The raw determining
slot is $b_j=\beta+b_{j,0}$ for a constant
integer vector. This convention matters: every later old polynomial
uses the actual integer slot, not an approximation to its value.
Structural numbers of path columns are fixed polynomials of the
input dimensions before any numerical precision is chosen.
Consequently all ambient dimensions along a branch are polynomial
in the current dimension forecasts. Sides are selected to support
the next layer and finally Lemma~\ref{ap:absolute}.

The return processes layers from $D$ down to $1$. Layers $j>s$
are passive: the new patch does not acquire their integer arguments.
At the final $s$ active layers its weights may increase by fixed
factors, while the old determining equations and their weights stay
exact. This distinction is why a fixed $D$ can be chosen before the
iteration, even though the number of slots grows.

The input certificate and the returned score carry different information.
Preparation needs only
the strict inequality in the input certificate: its selections do not charge
the inverse of that inequality's excess. Proposition~\ref{ap:scalar:forward}
then gives terminal means near the current threshold on a quantitatively
large set of paths. Lemma~\ref{ap:absolute} supplies the quantitative score
needed for ascent. Thus all subsequent approximation accuracies can be
chosen from the recorded budgets, independently of the original excess.

The analytic proofs are grouped separately.
Table~\ref{ap:interface-table}, at the entrance to that part,
locates each input statement and its proof.
